\documentclass[aspectratio=169,10pt]{beamer}
\usetheme{metropolis}
\usepackage[T1]{fontenc}
\usepackage[utf8]{inputenc}
\usepackage[ngerman]{babel}
\usepackage{booktabs}
\usepackage{amsmath}
\usepackage{graphicx}
\IfFileExists{results_path.tex}{\input{results_path.tex}}{\newcommand{\ResultsPath}{../results}}
\title{Neuronale Ersatzmodelle für Biegelinien}
\subtitle{Einfluss von Netzarchitektur, Datenmenge und Querschnittssprung}
\author{<Name>}
\date{\today}

\begin{document}
\maketitle

\begin{frame}{Fragestellung}
\begin{itemize}
  \item Wie genau approximieren Tanh-Netze analytische Biegelinien?
  \item Welchen Einfluss haben 80 und 160 Trainingsfälle?
  \item Wie verändern zwei versteckte Schichten und ein Querschnittssprung die Ergebnisse?
  \item Wie unterscheiden sich Interpolation und Extrapolation?
\end{itemize}
\end{frame}

\begin{frame}{Physikalische Referenz}
Für den Kragarm mit Endkraft gilt bei konstantem Querschnitt
\[
w(x)=\frac{F x^2(3L-x)}{6EI}.
\]
Beim Querschnittssprung integriert die Studie die Krümmung abschnittsweise. Durchbiegung und Neigung bleiben an der Sprungstelle stetig.
\end{frame}

\begin{frame}{Studiendesign}
\begin{center}
\begin{tabular}{lll}
\toprule
Geometrie & Netz & Trainingsfälle\\
\midrule
Konstant & 1 Layer, 1--15 Neuronen & 80, 160\\
Konstant & 2 Layer, 56 Kombinationen & 80, 160\\
Sprung & 1 Layer, 1--15 Neuronen & 80, 160\\
Sprung & 2 Layer, 56 Kombinationen & 80, 160\\
\bottomrule
\end{tabular}
\end{center}
Fünf Seeds je Architektur. Training, Validierung, Interpolation und Extrapolation verwenden getrennte Balkenfälle.
\end{frame}

\begin{frame}{Netz und Fehlermaß}
Das Netz approximiert die kraftnormierte Durchbiegung. Die Vorhersage lautet
\[
\widehat w = F\,w_N\,f_\theta(\text{Position, Geometrie}).
\]
Die Auswertung normiert den maximalen Fehler eines Falls mit der analytischen Enddurchbiegung $w(L)$.
\end{frame}

\begin{frame}{Konstanter Querschnitt: ein Layer, 80 Fälle}
\includegraphics[width=\linewidth,height=.76\textheight,keepaspectratio]{\ResultsPath/constant_1layer_80cases/main_plot.pdf}
\end{frame}

\begin{frame}{Konstanter Querschnitt: zwei Layer, 160 Fälle}
\includegraphics[width=\linewidth,height=.76\textheight,keepaspectratio]{\ResultsPath/constant_2layer_160cases/architecture_heatmaps.pdf}
\end{frame}

\begin{frame}{Querschnittssprung: ein Layer, 80 Fälle}
\includegraphics[width=\linewidth,height=.76\textheight,keepaspectratio]{\ResultsPath/step_1layer_80cases/main_plot.pdf}
\end{frame}

\begin{frame}{Querschnittssprung: zwei Layer, 160 Fälle}
\includegraphics[width=\linewidth,height=.76\textheight,keepaspectratio]{\ResultsPath/step_2layer_160cases/architecture_heatmaps.pdf}
\end{frame}

\begin{frame}{Vergleich der ausgewählten Modelle}
\includegraphics[width=\linewidth,height=.76\textheight,keepaspectratio]{\ResultsPath/overall_comparison.pdf}
\end{frame}

\begin{frame}{Ergebnisübersicht}
\begin{center}
\begin{tabular}{lrrrr}
\toprule
Geometrie & Layer & Fälle & Interpolation & Extrapolation\\
\midrule
\IfFileExists{study_summary.tex}{\input{study_summary.tex}}{Die Vollstudie erzeugt diese Tabelle.}\\
\bottomrule
\end{tabular}
\end{center}
\end{frame}

\begin{frame}{Einordnung}
\begin{itemize}
  \item Die Ergebnisse gelten für die gewählten Parameterbereiche und Trainingsbedingungen.
  \item Die Referenzdaten stammen aus analytischen Euler-Bernoulli-Lösungen.
  \item Die Kraftextrapolation folgt aus der eingebauten linearen Skalierung.
  \item Die Geometrieextrapolation prüft die Generalisierung des Netzes.
\end{itemize}
\end{frame}

\begin{frame}{Fazit}
Die Auswertung beantwortet die Frage nach einer geeigneten Architektur anhand von Validierungs-, Interpolations- und Extrapolationsfehlern. Die kleinste ausreichende Architektur wird über die Validierungsdaten ausgewählt.
\end{frame}
\end{document}
